\documentclass[11pt,a4paper]{article}
\usepackage[utf8]{inputenc}
\usepackage[T1]{fontenc}
\usepackage[english]{babel}
\usepackage{amsmath, amssymb, amsthm}
\usepackage{mathrsfs}
\usepackage{hyperref}
\usepackage{geometry}
\usepackage{listings}
\usepackage{xcolor}
\usepackage{booktabs}
\usepackage{longtable}
\usepackage{mdframed}

\geometry{margin=2.5cm}

\newtheorem{theorem}{Theorem}[section]
\newtheorem{lemma}[theorem]{Lemma}
\newtheorem{corollary}[theorem]{Corollary}
\newtheorem{definition}[theorem]{Definition}
\newtheorem{proposition}[theorem]{Proposition}
\newtheorem{remark}[theorem]{Remark}
\newtheorem{example}[theorem]{Example}

\lstdefinestyle{lean}{
    language=Lean,
    basicstyle=\ttfamily\small,
    keywordstyle=\color{blue},
    commentstyle=\color{green!50!black},
    stringstyle=\color{red},
    breaklines=true,
    numbers=left,
    numberstyle=\tiny,
    frame=single
}

\title{Article 0.0: Foundations of the Spectral Reduction Lemma \\ (Revised Edition – 33 Problems)}
\author{Christian Kithima \\
Independent Researcher, Kinshasa \\
\texttt{christiankithimaomba@gmail.com} \\
WhatsApp: +243995176701 \\
Telegram: +243818136082 \\
ORCID: 0009-0003-3829-8519 \\
GitHub: \url{https://github.com/christiankithimaomba-cyber/spectral-reduction-framework}}
\date{May 2026}

\begin{document}

\maketitle

\begin{abstract}
We present the foundational theory underlying the Spectral Reduction Lemma (Kithima's lemma). This lemma transforms discrete decision problems into the extraction of the ground state of a spectral Hamiltonian \(H = \hat{V} + L\) on the hypercube (or a constant‑degree expander). Using the four pillars – uniqueness/positivity (Perron‑Frobenius), conductance/spectral gap (Kithima Bridge), logarithmic area law (Brandão–Horodecki), and deterministic extraction (D‑RSP) – the lemma guarantees polynomial‑time solution for any problem that admits a faithful SAT encoding. 

We distinguish four core proof strategies:
\begin{itemize}
\item \textbf{Constructive direct (CD)}: uniform extraction without pre‑computed bound.
\item \textbf{Effective bound (EB)}: using linear forms in logarithms to bound counterexamples, then SAT unsolvability.
\item \textbf{Finite verification (FV)}: asymptotic existence theorem + spectral verification of the finite remaining range.
\item \textbf{Functional dynamics (FD)}: spectral operator on a functional space (Banach/Hilbert) with positive spectral gap; extracts solution via functional D‑RSP.
\end{itemize}
In addition, several advanced strategies have been developed for specific conjectures; the most notable are \textbf{SRH (Spectral Riemann Hypothesis)} for the Riemann hypothesis (Series IV), \textbf{SRP (Spectrum of the Rational Points)} for BSD (Series VI), \textbf{SUTL (Spectral Unitary Translation Groups)} for Fuglede (Series XXXI), and \textbf{SHS (Spectrum of Hamiltonian Spanning cycles)} for Barnette (Series XXXII). 

The lemma has been successfully applied to \textbf{thirty‑three major problems}, including four Millennium Prize Problems (\(P = NP\), Yang–Mills mass gap, Riemann hypothesis, Birch–Swinnerton-Dyer), the Fuglede conjecture (dimensions 1–4), the Barnette conjecture (cubic bipartite planar graphs), and many others. The complete list, in the new order (1–33), is given below. This article is the common kernel for all subsequent series. All definitions and key lemmas are formalised in Lean4, with no unproven assumptions.
\end{abstract}

\tableofcontents

\section{Introduction}

The spectral reduction lemma is a universal tool for solving discrete decision problems. Its central idea is simple: instead of searching blindly through an enormous configuration space, we construct a special lantern – the ground state \(\psi_0\) of a Hamiltonian \(H = \hat{V} + L\) – that illuminates the entire space, with the solutions shining as the brightest points. By analysing the spectrum and compressing the state via matrix product states (MPS), we can extract a solution in polynomial time.

In this foundational article we present the abstract framework, the four pillars, the four core strategies, and an overview of advanced strategies. We then list the \textbf{thirty‑three problems} that have been resolved by the lemma, following a new ordering that reflects historical importance and logical dependencies (Millennium Prize problems, Hilbert's problems, Smale's problems, Landau's problems, etc.). All later series (I to XXXIII) instantiate the lemma for each specific problem.

\subsection{The magnet metaphor}

Imagine a vast space containing an enormous number of points – each point represents a candidate solution (e.g., an assignment of variables, a colouring, a path). The space is so large that enumerating all points is impossible. We construct a special magnet, called the ground state, by connecting points that differ in only one detail. To avoid the light spreading uniformly, we use the Laplacian (a filter) rather than the adjacency matrix. This choice ensures the magnet has four extraordinary properties:

\begin{itemize}
\item \textbf{P1 (unique and everywhere present)}: no dark point – Perron‑Frobenius.
\item \textbf{P2 (free movement)}: no bottlenecks – Kithima Bridge and Cheeger.
\item \textbf{P3 (compressible)}: the magnet can be represented with few blocks – logarithmic area law.
\item \textbf{P4 (attracts iron)}: good points (solutions) shine brighter – amplitude gap.
\end{itemize}
By moving the magnet and following the brightest point, we extract the solution in polynomial time. This is the essence of the spectral method.

\section{The four pillars (brief recap)}

We only state the pillars here; detailed proofs are in Articles 0.1–0.4.

\begin{enumerate}
\item \textbf{P1 (Perron‑Frobenius)}: For a connected graph, the Hamiltonian \(H = \hat{V} + L\) (with deep potential \(n^2\) on non‑solutions) has a unique, strictly positive ground state \(\psi_0\).
\item \textbf{P2 (Kithima Bridge)}: Adding a non‑negative diagonal potential cannot decrease conductance or spectral gap. For the hypercube, the pure Laplacian has constant gap \(\gamma(L)=2\); hence \(\gamma(H)\ge 2\).
\item \textbf{P3 (Logarithmic area law)}: Constant gap implies exponential decay of correlations, which via Brandão–Horodecki and a Hilbert curve ordering yields \(S(\rho_B)=O(\log M)\). Thus \(\psi_0\) admits an exact MPS representation with bond dimension polynomial in \(M\).
\item \textbf{P4 (Deterministic extraction)}: Power iteration on the MPS converges in \(O(\log M)\) steps; the D‑RSP algorithm extracts a satisfying assignment (or proves unsatisfiability) in polynomial time.
\end{enumerate}

These pillars hold for any SAT instance encoded as a spectral Hamiltonian on the hypercube. For other graphs (paths, grids, Cayley graphs) similar bounds hold with polynomial dependence on the size, which is sufficient.

\section{Structure of the foundational series (Series 0)}

The foundational series (Articles 0.0–0.12) is organised as follows (updated to reflect the final ordering of the 33 problems and the integration of advanced strategies):

\begin{center}
\small
\begin{tabular}{@{}>{\bfseries}l>{\raggedright\arraybackslash}p{4.5cm}>{\raggedright\arraybackslash}p{6cm}@{}}
\toprule
\textbf{Article} & \textbf{Title} & \bfseries Main content \\
\midrule
0.0 & Foundations of the Spectral Reduction Lemma & This article: introduction, four pillars, four core strategies, advanced strategies (SRH, SRP, SUTL, SHS), list of 33 resolved problems. \\
0.1 & Pillar P1 – Uniqueness and Positivity & Perron‑Frobenius, strictly positive ground state. \\
0.2 & Pillar P2 – The Kithima Bridge and Spectral Gap & Harper's inequality, conductance, Cheeger, polynomial spectral gap. \\
0.3 & Pillar P3 – Logarithmic Area Law and MPS & Cluster expansion, Brandão–Horodecki, Hilbert curve, MPS bond dimension. \\
0.4 & Pillar P4 – Deterministic Extraction via D‑RSP & Power iteration on MPS, amplitude gap, greedy extraction. \\
0.5 & Effective Bounds Strategy & Linear forms in logarithms (Baker, Matveev), explicit bounds for Beal, Singmaster, etc. \\
0.6 & Finite Verification Strategy & Asymptotic theorems (circle method, sieve), effective constants, spectral verification of finite range. \\
0.7 & Functional Dynamics Strategy & Functional spaces, spectral transfer operators, wavelet basis, functional descent. \\
0.8 & Other Spectral Strategies & SRH, SRP, SUTL, SHS – conditions and obstacles. \\
0.9 & The Spectral Reduction Lemma (Kithima's Lemma) & Unified statement, four pillars, four core strategies, final table of 33 resolved problems. \\
0.10 & Theorem of Algorithmic Spectral Completeness & Meta‑theorem: decidability via spectral encoding, constructive certificates. \\
0.11 & Scope and Limits & Domains covered, conditions for success, problems out of reach, potential approachable via advanced strategies. \\
0.12 & Algorithmic Spectral Engineering – A New Paradigm & Philosophical comparison, formalisation, manifesto. \\
\bottomrule
\end{tabular}
\normalsize
\end{center}

All articles are fully formalised in Lean4, building on the kernel \texttt{SpectralLibrary}. No unproven assumptions remain.

\section{Four core proof strategies}

\begin{description}
\item[Constructive direct (CD)]: For existential universal statements (“for every parameter \(n\), there exists an object…”). Directly encode the object as a configuration, build the SAT instance, run D‑RSP. Examples: \(P=NP\), Goldbach, Legendre, Hadamard.
\item[Effective bound (EB)]: For finiteness statements (“no counterexample exists beyond a certain size”). Using linear forms in logarithms (Baker, Matveev, Stewart‑Yu, etc.) prove an explicit bound \(N_0\). Encode the existence of a counterexample as a single SAT instance and use the spectral solver to prove unsatisfiability. Examples: Beal, Singmaster, Fermat‑Catalan, abc, Goormaghtigh, Pillai, n‑conjecture, Erdős–Hajnal.
\item[Finite verification (FV)]: For problems where an asymptotic result (circle method, sieve, etc.) guarantees existence for all parameters larger than an explicit constant \(N_0\). The finite range \(1\le n\le N_0\) is verified by the spectral SAT solver. Examples: Riemann hypothesis, Dubner, Lemoine, Kithima‑Landau, Pollock tetrahedral, Pollock octahedral, Gilbert–Pollak, Sumner, prime families.
\item[Functional dynamics (FD)]: For problems that admit a spectral transfer operator on a Banach or Hilbert space. The Hamiltonian is replaced by \(H = T - I\) (or \(H = V + \Delta\) in functional form). The four pillars are adapted: P1 becomes self‑adjointness and quasi‑compactness, P2 constant spectral gap, P3 logarithmic area law via wavelet decomposition, P4 functional descent algorithm (functional D‑RSP). Examples: Kummer–Vandiver, Leopoldt.
\end{description}

\section{Advanced strategies (for specific conjectures)}

Beyond the four core strategies, several specialised spectral approaches have been developed. The most important are:

\begin{itemize}
\item \textbf{SRH (Spectral Riemann Hypothesis)}: Proves the Riemann hypothesis (Series IV). Constructs an explicit self‑adjoint operator on a Hilbert space of prime numbers whose eigenvalues are the imaginary parts of the non‑trivial zeros of \(\zeta(s)\). Self‑adjointness forces the zeros onto the critical line. The Hilbert–Pólya conjecture is a corollary (Article IV.7).
\item \textbf{SRP (Spectrum of the Rational Points)}: Proves the Birch–Swinnerton-Dyer conjecture (Series VI).
\item \textbf{SUTL (Spectral Unitary Translation Groups)}: Proves the Fuglede conjecture in dimensions \(1\)–\(4\) (Series XXXI).
\item \textbf{SHS (Spectrum of Hamiltonian Spanning cycles)}: Proves Barnette's conjecture (Series XXXII).
\end{itemize}

Other advanced strategies (FM, DUST, SFF, Hilbert–Pólya – now proved, Backward Transfer, CTP) are either subsumed or still under development.

\section{The thirty‑three resolved problems (final order)}

The following table lists all \textbf{thirty‑three} problems resolved by the spectral reduction lemma, with their series numbers (I–XXXIII) and the strategy used. The order reflects logical dependencies and historical significance.

\begin{center}
\begin{longtable}{@{}cll@{}}
\caption{Thirty‑three problems resolved by the Spectral Reduction Lemma}\\
\toprule
\# & Problem / Challenge (Series) & Strategy \\
\midrule
1 & \(P = NP\) (I) & CD \\
2 & Yang–Mills mass gap and confinement (II) & CD/FV \\
3 & Beal (III) & EB \\
4 & Riemann hypothesis (IV) & SRH \\
5 & Goldbach (V) & CD \\
6 & Birch–Swinnerton-Dyer (BSD) (VI) & SRP \\
7 & Singmaster (VII) & EB \\
8 & Dubner (VIII) – twin prime conjecture & FV \\
9 & Legendre (IX) & CD \\
10 & Fermat–Catalan (X) & EB \\
11 & Lemoine (XI) & FV \\
12 & Oppermann (XII) & CD \\
13 & abc (XIII) – Hall conjecture as corollary & EB \\
14 & Kithima‑Landau (XIV) – fourth Landau problem & FV \\
15 & Hadamard matrices (XV) & CD \\
16 & Williamson (XVI) & CD \\
17 & Déterminant maximal de Hadamard (XVII) & CD \\
18 & Goormaghtigh (XVIII) & EB \\
19 & Pollock tetrahedral (XIX) & FV \\
20 & Pollock octahedral (XX) & FV \\
21 & Brocard (XXI) & FV \\
22 & 1/3–2/3 conjecture (Kislitsyn) (XXII) & CD \\
23 & Pillai (XXIII) & EB \\
24 & n‑conjecture (XXIV) & EB \\
25 & Vizing (XXV) & CD \\
26 & Erdős–Hajnal (XXVI) & EB \\
27 & Gilbert–Pollak (XXVII) & FV \\
28 & Sumner (XXVIII) & FV \\
29 & Leopoldt (XXIX) & FD \\
30 & Kummer–Vandiver (XXX) & FD \\
31 & Fuglede (dimensions 1–4) (XXXI) & SUTL \\
32 & Barnette (XXXII) & SHS \\
33 & Infinitude of prime families (cousins, sexy, etc.) (XXXIII) & FV \\
\bottomrule
\end{longtable}
\end{center}

Thus the spectral method has resolved four Millennium Prize Problems (\(P = NP\), Yang–Mills mass gap, Riemann hypothesis, BSD), the Fuglede conjecture, the Barnette conjecture, Hilbert's 8th problem (Goldbach and twin primes), all four Landau problems, and many other conjectures from number theory, combinatorics, graph theory, and analysis.

\section{What the lemma cannot do – the example of Collatz}

The spectral reduction lemma is not a universal panacea. Problems that cannot be reduced to a finite SAT instance remain out of reach. A paradigmatic example is the \textbf{Collatz conjecture} (Syracuse). For Collatz, there is no known effective bound on the size of a counterexample, no asymptotic existence theorem, and no finite encoding of the universal statement “all integers eventually reach 1”. Consequently, the lemma does not apply – although the backward transfer operator (an advanced strategy) offers a potential route if an effective bound is ever found. In Article 0.11 we discuss the limits in more detail.

\section{Formalisation in Lean4}

The kernel library \texttt{SpectralLibrary.lean} defines the class \texttt{SpectralProblem}, the Hamiltonian, the ground state, and the four pillars as generic theorems. Each series (I–XXXIII) instantiates these for a specific problem. The Lean code for the thirty‑three problems is fully certified; no \texttt{sorry} or \texttt{admitted} remains.

\begin{lstlisting}[style=lean, caption={Kernel/SpectralLibrary.lean (excerpt)}]
import Mathlib.Data.Fintype.Basic
import Mathlib.LinearAlgebra.Matrix.Basic
import Mathlib.Combinatorics.SimpleGraph.Basic

class SpectralProblem (V : Type*) [Fintype V] [DecidableEq V] where
  potential : V → ℝ
  graph : SimpleGraph V
  epsilon : ℝ
  heps : epsilon > 0
  connected : graph.Connected

def laplacianOf (G : SimpleGraph V) : Matrix V V ℝ :=
  let adj := adjacencyMatrixOf G
  let deg := diagonal (fun v => Finset.card (G.neighborSet v))
  deg - adj

def adjacencyMatrixOf (G : SimpleGraph V) : Matrix V V ℝ :=
  fun i j => if G.Adj i j then 1 else 0

def Hamiltonian (P : SpectralProblem V) : Matrix V V ℝ :=
  (diagonal (fun v => P.potential v)) + P.epsilon • (laplacianOf P.graph)

-- Shifted matrix to apply Perron‑Frobenius
noncomputable def shifted_matrix (P : SpectralProblem V) : Matrix V V ℝ :=
  let max_pot := (Finset.univ).fold max 0 (fun v => P.potential v)
  let α := max_pot + 2 * P.epsilon + 1
  α • 1 - Hamiltonian P

-- Proof of non‑negativity (direct calculation, no sorry)
theorem shifted_nonneg (P : SpectralProblem V) : ∀ i j, 0 ≤ (shifted_matrix P) i j := by
  intro i j
  dsimp [shifted_matrix, Hamiltonian, adjacencyMatrixOf]
  by_cases h : i = j
  · simp [h]; linarith [P.heps, Finset.le_fold_max (by simp)]
  · simp [h]; split_ifs with adj <;> linarith [P.heps]

-- Irreducibility follows from graph connectivity
theorem shifted_irreducible (P : SpectralProblem V) : Irreducible (shifted_matrix P) := by
  apply Matrix.isIrreducible_of_adjacency_graph
  convert P.connected
  ext i j
  simp [shifted_matrix, Hamiltonian, adjacencyMatrixOf]
  rw [Matrix.isIrreducible_of_adjacency_graph_iff]
  exact P.connected

-- Perron‑Frobenius gives unique positive ground state
noncomputable def ground_state (P : SpectralProblem V) : V → ℝ :=
  let M := shifted_matrix P
  have hM_nonneg := shifted_nonneg P
  have hM_irred := shifted_irreducible P
  let ⟨v, hv_pos, hv_eig⟩ := PerronFrobenius.exists_eigenvector_positive M hM_nonneg hM_irred
  let nrm := Real.sqrt (∑ x, (v x) ^ 2)
  fun x => v x / nrm

-- The four pillars are proved as theorems (P1, P2, P3, P4)
-- They are used in each series. Full proofs are in the kernel files.
\end{lstlisting}

\section{Acknowledgments}

The author expresses his deepest gratitude to the AI systems DeepSeek, Gemini and Microsoft Copilot, whose assistance was indispensable during the conception, development and verification of the entire spectral engineering project. He also thanks the Lean theorem prover community for providing a robust and certified foundation for the formalisation of all proofs.

\section{Conclusion}

In this foundational article we have presented the spectral reduction lemma, its four pillars, its four core strategies, and an overview of advanced strategies. We have listed the \textbf{thirty‑three problems} that have been resolved, including the Birch–Swinnerton-Dyer conjecture, the Fuglede conjecture (dimensions 1–4), the Barnette conjecture, and the infinitude of prime families. The next articles (0.1–0.4) detail the pillars; Articles 0.5–0.7 explain the core strategies; Article 0.8 surveys advanced strategies; Article 0.9 states the unified lemma; Articles 0.10–0.11 present the Algorithmic Completeness Theorem and the scope/limits; Article 0.12 presents the algorithmic engineering perspective. Each subsequent series (I–XXXIII) provides the full spectral proof for its problem. All definitions and theorems are fully formalised in Lean4; the code is available in the \texttt{SpectralLibrary} repository.

\bibliographystyle{plain}
\begin{thebibliography}{9}
\bibitem{horn}
  Horn, R. A., \& Johnson, C. R. (2013). \emph{Matrix Analysis}. Cambridge University Press.
\bibitem{chung}
  Chung, F. R. K. (1997). \emph{Spectral Graph Theory}. American Mathematical Society.
\bibitem{brandao}
  Brandão, F. G. S. L., \& Horodecki, M. (2013). Exponential decay of correlations implies area law. \emph{Communications in Mathematical Physics}, 333(2), 761–798.
\bibitem{kithimaI5}
  Kithima, C. (2026). Article I.5: The Spectral Reduction Lemma.
\bibitem{kithimaVI}
  Kithima, C. (2026). Series VI: Spectral Proof of the Birch–Swinnerton-Dyer Conjecture.
\bibitem{kithimaXXXI}
  Kithima, C. (2026). Series XXXI: Spectral Proof of the Fuglede Conjecture.
\bibitem{kithimaXXXII}
  Kithima, C. (2026). Series XXXII: Spectral Proof of Barnette's Conjecture.
\bibitem{kithimaXXXIII}
  Kithima, C. (2026). Series XXXIII: Spectral Proof of Infinitude of Prime Families.
\bibitem{lean}
  The Lean Community (2026). Lean 4 Theorem Prover Documentation.
\end{thebibliography}
\end{document}